% Source: https://artofproblemsolving.com/community/c3h1414339_geometry
% Stable identity: aops-topic-1414339-q01
The semicircle has diameter $XW=UV=20$, so its radius is $10$ and its highest point $Z$ is $10$ above the midpoint $M$ of $XW$.

Use coordinates with the floor as the $xy$-plane: $U=(0,0,0)$, $V=(20,0,0)$, $W=(20,30,0)$, $X=(0,30,0)$. Then $M=(10,30,0)$ and $Z=(10,30,10)$.

\begin{asy}
usepackage("amsmath");
size(6.5cm);
defaultpen(fontsize(10pt));
pair P(real x,real y,real z){return (x,z)+y*0.45*dir(35);}
pair U=P(0,0,0),V=P(20,0,0),Wp=P(20,30,0),X=P(0,30,0),Z=P(10,30,10),Mx=P(10,30,0);
import graph;
path wall=graph(new pair(real t){return P(10+10*cos(t),30,10*sin(t));},0,pi);
fill(wall--cycle,lightgray);
draw(wall);
draw(U--V--Wp--X--cycle,linewidth(0.9));
draw(Mx--Z,gray+dashed);
draw(U--Z,red+linewidth(0.9)); draw(V--Z,red+linewidth(0.9));
dot(Z,red); dot(Mx);
label("$U$",U,SW); label("$V$",V,SE); label("$W$",Wp,E); label("$X$",X,SE); label("$Z$",Z,N,red); label("$M$",Mx,SE);
label("$20$",(U+V)/2,S); label("$30$",(V+Wp)/2,SE); label("$10$",(Mx+Z)/2,E,gray);
label("\mbox{}",(min(currentpicture,user=true)+max(currentpicture,user=true))/2,black); // keeps the text after the figure black when LaTeX typesets the labels inline
\end{asy}

So
\[
UZ=\sqrt{10^2+30^2+10^2}=\sqrt{1100}=10\sqrt{11},
\]
and by symmetry $VZ=10\sqrt{11}$ as well. The perimeter is
\[
20+20\sqrt{11}\approx20+66.33=86.33,
\]
which is closest to $\boxed{86}$.
